\chapter{Conjunctivitis}

\section*{Definition and Overview}
Conjunctivitis, commonly known as ``pink eye'', is inflammation of the conjunctiva, the thin, transparent membrane that lines the inner surface of the eyelids and covers the white part of the eye. It is one of the most frequently encountered eye conditions in primary care and is usually mild and self-limiting. The principal causes are viral infection, bacterial infection, and allergy, although chemical irritants, foreign bodies, and contact lens-related problems can also produce the same clinical picture. Most cases resolve within one to two weeks, but the appropriate treatment and the degree of urgency depend heavily on identifying the underlying cause. Recognition of the features that signal a more serious condition, such as pain, photophobia, or reduced vision, is essential, because these point away from simple conjunctivitis and towards conditions that threaten the eye.

\section*{Epidemiology}
Conjunctivitis is extremely common across all age groups and accounts for a substantial proportion of presentations to general practitioners, emergency departments, and community pharmacies. Viral conjunctivitis is the most frequent form in adults and is highly contagious, spreading readily in schools, childcare centres, households, and workplaces through hand-to-eye contact and shared objects. Bacterial conjunctivitis is proportionally more common in children than in adults and is a frequent cause of school and childcare exclusion. Allergic conjunctivitis affects a significant fraction of the population at some point in life, with higher prevalence among people who also have hay fever, asthma, or eczema, and it shows seasonal peaks that mirror pollen and mould seasons. Neonatal conjunctivitis (ophthalmia neonatorum) is relatively uncommon in high-income settings but remains important because of the risk of sight-threatening and systemic complications. Both viral and bacterial forms cluster in the winter months, while allergic forms peak in spring and summer.

\section*{Aetiology and Pathophysiology}
\medskip
\noindent The causes of conjunctivitis can be grouped as follows:
\begin{itemize}
    \item \textbf{Viral:} adenoviruses are the most common cause, with enteroviruses, herpes simplex virus, and other respiratory viruses implicated less often; infection is highly contagious and may accompany an upper respiratory tract infection
    \item \textbf{Bacterial:} \textit{Staphylococcus aureus}, \textit{Streptococcus pneumoniae}, \textit{Haemophilus influenzae}, and \textit{Moraxella catarrhalis} are typical isolates; children and contact lens wearers are most commonly affected
    \item \textbf{Allergic:} an immunoglobulin E-mediated hypersensitivity reaction to pollens, dust mites, pet dander, or mould triggers mast-cell degranulation and histamine release within the conjunctiva
    \item \textbf{Irritant:} smoke, chlorine, chemical splashes, air pollution, and preservatives in eye drops cause a non-specific inflammatory response
    \item \textbf{Contact lens-related:} poor lens hygiene, overwear, or solution reactions, with a risk of secondary bacterial or fungal keratitis
    \item \textbf{Neonatal:} infection acquired during passage through the birth canal, most importantly chlamydia or gonorrhoea, requiring urgent specialist review
    \item \textbf{Inflammatory:} conditions such as dry eye, blepharitis, and systemic autoimmune disease can produce chronic conjunctival inflammation
\end{itemize}

Whatever the trigger, the conjunctival vasculature dilates and becomes more permeable, allowing fluid and inflammatory cells to leak into the tissue. This produces the cardinal features of redness, swelling, and discharge. In viral and allergic disease the discharge tends to be watery, whereas in bacterial disease it is typically mucopurulent.

\section*{Clinical Features}
The features vary somewhat by cause, but typical symptoms include:
\begin{itemize}
    \item Redness and irritation of the white of the eye, usually starting in one eye and sometimes spreading to both
    \item A gritty or burning sensation, as if there is sand in the eye
    \item Watery discharge in viral and allergic conjunctivitis
    \item Thick, yellow-green, or crusting discharge in bacterial conjunctivitis, often with the eyelids stuck together on waking
    \item Itching, which is prominent in allergic conjunctivitis and typically affects both eyes
    \item Swelling of the eyelids and conjunctiva (chemosis)
    \item Mild foreign-body sensation; vision is normally preserved, and blurred or reduced vision is not typical and warrants review
\end{itemize}
The presence of significant eye pain, photophobia, or a reduction in visual acuity suggests corneal involvement or intraocular inflammation rather than simple conjunctivitis and should be taken seriously.

\section*{Differential Diagnosis}
\medskip
\noindent A number of conditions produce a red eye and may be mistaken for conjunctivitis:
\begin{itemize}
    \item \textbf{Subconjunctival haemorrhage:} a painless, well-defined red patch caused by a broken blood vessel, with no discharge
    \item \textbf{Keratitis:} corneal inflammation causing pain, photophobia, and reduced vision, which can be sight-threatening
    \item \textbf{Anterior uveitis (iritis):} inflammation inside the eye, typically painful with a red ring around the cornea
    \item \textbf{Acute angle-closure glaucoma:} an emergency with severe pain, headache, blurred vision, and nausea
    \item \textbf{Foreign body or corneal abrasion:} localised pain and a history of something entering the eye
    \item \textbf{Dry eye syndrome:} chronic irritation and redness with a grittiness that worsens through the day
    \item \textbf{Blepharitis:} inflammation of the eyelid margins causing burning, crusting, and redness
    \item \textbf{Episcleritis and scleritis:} deeper, more localised inflammation that is often more painful than conjunctivitis
    \item \textbf{Contact lens-related keratitis:} potentially serious, requiring urgent assessment
\end{itemize}

\section*{When to Seek Medical Care}
Most cases of conjunctivitis settle on their own, but medical assessment is advisable if symptoms are severe, fail to improve within a few days, the affected person wears contact lenses, or the person is a newborn or young infant.

\textbf{Call 000 (or local emergency services) for:}
\begin{itemize}
    \item Significant eye pain or marked sensitivity to light
    \item Blurred or reduced vision
    \item Rapidly worsening redness, especially involving the cornea
    \item Symptoms in a newborn baby
    \item A chemical splash or suspected foreign body in the eye
    \item Redness associated with fever, nausea, or headache suggesting systemic or intraocular disease
\end{itemize}

\section*{Diagnosis and Investigations}
\medskip
\noindent In most cases the diagnosis is clinical and no investigation is required.
\begin{itemize}
    \item \textbf{History:} laterality, duration, type of discharge, contact with someone with a red eye, associated cold symptoms, allergies, contact lens use, trauma, and occupational or chemical exposures
    \item \textbf{Examination:} inspection of the conjunctiva, eyelids, and cornea with a torch or slit lamp; assessment of the pupils; and examination for pre-auricular lymphadenopathy
    \item \textbf{Visual acuity testing:} to confirm that vision is unaffected
    \item \textbf{Fluorescein staining:} to detect corneal abrasions or ulcers
    \item \textbf{Microbial swabs:} reserved for severe, recurrent, or suspected gonococcal or chlamydial infection, for newborns, and for contact lens-related cases
    \item \textbf{Referral:} to an ophthalmologist when vision is affected, the diagnosis is uncertain, symptoms worsen despite treatment, or there is any suspicion of keratitis, uveitis, or angle-closure glaucoma
\end{itemize}

\section*{Management}
\medskip
\noindent Management depends on the cause, and many cases require no treatment at all.
\begin{itemize}
    \item \textbf{Viral conjunctivitis:} self-limiting; cool compresses, preservative-free artificial tears, and simple analgesia relieve symptoms; resolution usually occurs within one to two weeks
    \item \textbf{Bacterial conjunctivitis:} frequently clears without antibiotics; if prescribed, topical antibiotic drops are used for a short course in severe, persistent, or high-risk cases
    \item \textbf{Allergic conjunctivitis:} avoidance of the trigger is central; topical antihistamine or mast-cell stabiliser drops and cool compresses provide symptomatic relief, and oral antihistamines may help in seasonal disease
    \item \textbf{General measures:} gently clean the eyelids with cooled boiled water and a clean cloth; stop wearing contact lenses until the eye is fully settled
    \item \textbf{Hygiene:} wash hands frequently, avoid rubbing the eyes, and do not share towels, face cloths, pillows, or eye cosmetics while contagious
    \item \textbf{Follow-up:} review if there is no improvement within several days or if symptoms worsen, particularly if pain or blurred vision develops
\end{itemize}

\section*{Complications}
\medskip
\noindent Although most cases are uncomplicated, the following can occur:
\begin{itemize}
    \item Spread of infection to household members, classmates, and colleagues
    \item Corneal involvement (keratitis) in some viral, bacterial, and contact lens-related infections, which can threaten vision
    \item Chronic or recurrent conjunctival inflammation with, rarely, scarring of the conjunctiva
    \item Secondary bacterial infection on top of a viral conjunctivitis
    \item Harm from self-treatment, including eye damage from contaminated, expired, or non-sterile drops
    \item Serious complications of neonatal conjunctivitis, including corneal scarring and systemic sepsis, when untreated
\end{itemize}

\section*{Prognosis and Outcomes}
The outlook is generally excellent. Viral conjunctivitis usually resolves within one to two weeks without treatment, bacterial conjunctivitis typically clears within a few days of appropriate therapy, and allergic conjunctivitis responds well to trigger avoidance and topical treatment. Permanent visual loss is extremely rare in straightforward cases. People with recurrent allergic conjunctivitis may require treatment through each allergy season, but most manage their symptoms effectively with preventive measures and medication. Patients who wear contact lenses and those with severe or recurrent disease should be counselled to seek review early, because the consequences of untreated corneal involvement are far greater.

\section*{Prevention and Self-Care}
\begin{itemize}
    \item Wash hands frequently, especially after touching the eyes or face
    \item Avoid touching or rubbing the eyes
    \item Do not share towels, face cloths, pillows, or eye cosmetics
    \item Discard eye makeup and replace contact lens solution and cases if an infection occurs
    \item If you wear contact lenses, follow cleaning and replacement advice and do not wear lenses while the eye is infected
    \item For allergic conjunctivitis, reduce exposure to known triggers and consider allergy-proofing the home
    \item Stay away from school, childcare, or work while the discharge is heavy, if advised by the treating team
    \item Use separate hand towels and dispose of used tissues carefully during an infectious episode
\end{itemize}

\section*{Sources and Further Reading}
The following organisations provide reliable, current information on conjunctivitis and red eye:
\begin{itemize}
    \item NHS. \textit{Health information on conjunctivitis and eye conditions.} https://www.nhs.uk/conditions/
    \item Mayo Clinic. \textit{Patient education resources on eye conditions.} https://www.mayoclinic.org/diseases-conditions
    \item Centers for Disease Control and Prevention (CDC). \textit{Public health information on conjunctivitis (pink eye).} https://www.cdc.gov/conjunctivitis/
    \item MedlinePlus. \textit{Consumer health information on eye disorders.} https://medlineplus.gov/
    \item MSD Manual. \textit{Professional and consumer reference on eye disorders.} https://www.msdmanuals.com/home/eye-disorders
\end{itemize}
